%---------------------------------------------------------------------
%	Abstract
%---------------------------------------------------------------------
\chapter{Abstract}				
\label{ch:abs}
%---------------------------------------------------------------------
% TAKEN FROM UCT WEBSITE
% The Doctoral Degrees Board recommends that candidates include an abstract 
% that fits onto one page, and which includes the author's full name, thesis 
% title and date. The text should not exceed 350 words. Candidates may also 
% include a more substantial summary of their work, in addition to the abstract.
% Follow the same format for an MSc. abstract.

% The abstract should stand on its own. It should answer the following questions:
%	1.	What did the author do? What ideas, notions, hypotheses, 
%		concepts, theories or thoughts were investigated?
%	2.	How did the author do the work? What data were generated and used? 
%		What was the origin of the data? How were data gathered? 
%		What tests, scales, indices, or summary measures were used? 
%		In other words, how was the analysis and/or synthesis done?
%	3.	What were the conclusions and what were the significant findings?
%
% Some studies cannot readily be summarised in this way and require more 
% descriptive abstracts. Do not use telegraphic phrases. Do not repeat 
% information given in the title. Do not use abbreviations. The purpose of 
% an abstract is to enable a researcher/examiner to understand the essential 
% hypothesis, method and findings of the research.
%---------------------------------------------------------------------
\begin{center}
	\textbf{\Large \titl}\\
			\vskip 0.2cm
			\auth\\
			\vskip 0.2cm
	\textit{\footnotesize \FullDate}
			\vskip 1cm
\end{center}
%---------------------------------------------------------------------+

Sea ice extent and floe morphology in the polar regions play a critical role in modulating ice dynamics and air-sea heat exchange, yet floe-scale properties remain poorly constrained in Earth System Models. Improved characterization of floe properties, and their interaction with their environment, specifically their interactions with ocean waves, is essential for enhancing these models, particularly those that simulate polar climate processes and the cryosphere's energy budget.

This study makes use of two LiDAR datasets to investigate the estimation of key physical properties of sea ice floes (size and shape), as well as surface-wave parameters (amplitude and frequency) in young ice fields in the Antarctic Marginal Ice Zone, focusing on the Marginal Ice Zone in winter, when the sea ice extent reaches its maximum. The first dataset is derived from controlled experiments conducted in the Aalto University Wave and Ice Tank, while the second dataset comprises field measurements collected in the Antarctic Marginal Ice Zone during the SCALE Winter 2022 expedition aboard the South African research vessel \textit{S.A. Agulhas II}. The measurements were obtained using a Livox Avia LiDAR sensor integrated with an Inertial Measurement Unit, enabling high-resolution point cloud acquisition and motion tracking in shipborne conditions. 

Data processing and analysis is conducted in MATLAB through a custom pipeline consisting of sensor characterization, pre-processing and feature extraction stages. The sensor characterization stage involves assessing the LiDAR sensor's performance in terms of scene reconstruction accuracy and behaviour under simulated data capture conditions. Pre-processing steps include motion compensation using Kalman filtering informed by Inertial Measurement Unit data, as well as point cloud filtering to mitigate noise from environmental artifacts such as snowfall. Feature extraction focuses on estimating wave parameters through plane fitting techniques and clustering points to identify individual floes using density-based clustering algorithms.

The results obtained demonstrate that despite challenges posed by the dynamic and harsh Antarctic environment and limitations of the specific LiDAR sensor used to collect data, the developed processing pipeline effectively extracts meaningful information regarding sea ice floe characteristics and wave dynamics. LiDAR-derived measurements of simulated wave parameters show good agreement with ground truth data, validating the approach. The clustering algorithms showed promise in isolating individual floes, although further refinement is needed to enhance accuracy in complex ice conditions. These findings demonstrate the potential of LiDAR-based observation systems to provide floe-scale observations that can better inform physics for polar climate modelling.

%---------------------------------------------------------------------
